\section{Introduction}\label{sec:intro}

Quantum error correction protects logical information by measuring stabilizers of a code repeatedly and decoding the
resulting detector events. The surface code~\cite{kitaev2003anyons,dennis2002topological,fowler2012surface} is the
leading candidate, and recent experiments have shown logical error rates that fall as the code distance
grows~\cite{google2023scaling,google2025below}. Designing such experiments, estimating thresholds and benchmarking
decoders~\cite{higgott2022pymatching,higgott2025sparse,delfosse2021unionfind} all depend on Monte Carlo sampling:
a noisy stabilizer circuit is sampled millions to billions of times, and every sample (a shot) yields one bit per
detector and per logical observable. The statistical error of a logical error rate $p_L$ estimated from $N$ shots is
about $\sqrt{p_L/N}$, so resolving the small error rates that matter for fault tolerance needs very many shots.

Stim~\cite{gidney2021stim} made this workload fast on CPUs. It computes one noiseless reference sample of the circuit
with a stabilizer tableau~\cite{gottesman1998heisenberg,aaronson2004chp} and then samples noisy shots by propagating
only the Pauli frame~\cite{knill2005noisy}, the deviation of each shot from the reference, through the circuit. A
frame update is a few bitwise operations per gate, applied to many shots packed into one machine word. The work is
therefore regular, branch free and independent between shots: an obvious match for a deeply pipelined hardware
datapath with wide words, and also for many devices working side by side.

This paper treats that observation as a systems question for FPGA accelerators. In earlier work we built a
distributed statevector simulator on four Alveo U55C cards in one host, in which every card is a node that holds a
slice of the state, and amplitudes move between nodes through host memory~\cite{ali2026multifpga}. Treating the cards
as the nodes of a small cluster, with the host acting as interconnect and scheduler, is the model that let four cards
work as one machine there, and it is the model we adopt here.
Pauli frame sampling fits the same node model, with one important difference: shots are independent, so nodes never
need to exchange data. The questions are then how fast one node can run, whether four nodes scale without
interfering through the shared host, and how the result compares with Stim running on all cores of the same host.

A first single card kernel for this workload, which we built earlier in this project and which has not been
published (Section~\ref{sec:first}), validated the approach end to end on hardware but was slower than one Stim process on every circuit tested, because its instruction
loop was not pipelined and each launch sampled only 64 shots. Here we redesign the kernel around the three costs that
dominated it and run it on four nodes. The contributions are:
\begin{itemize}
\item A dataflow kernel that processes one instruction slot per clock cycle for \NumLanes{} shots
(Section~\ref{sec:arch}): a pipelined counter based noise stage, a Pauli frame in a one read, one write memory that
completes every slot two cycles after reading it, and a detector fold stage that works on a per batch measurement
record while the next batch is simulated.
\item A compiler schedule that removes all frame hazards ahead of time: two slots on the same qubit are placed at
least \NumHazardDistance{} slots apart, so the frame pipeline needs no hazard logic and still runs at the full rate
(Section~\ref{sec:frame}). On a distance 5 surface code the schedule adds only its \NumNopsSurfFive{} padding
slots.
\item A node model for independent sampling on several cards, in which every node, kernel instance and batch draws
its own Threefry stream~\cite{salmon2011random123} and the host only launches work and collects results
(Section~\ref{sec:nodes}).
\item A verification chain from a software model through C simulation and RTL cosimulation to every card, closed by
statistical comparison with Stim and by decoded logical error rates up to distance 7
(Sections~\ref{sec:verification} and~\ref{sec:res-correct}).
\item Measurements on one, two and four cards against Stim with one process and with all cores of the same host
(Section~\ref{sec:res-throughput}).
\end{itemize}

Section~\ref{sec:related} reviews related work, Section~\ref{sec:nodes} describes the node model and
Section~\ref{sec:arch} the kernel. Section~\ref{sec:verification} gives the verification method,
Section~\ref{sec:setup} the experimental setup and Section~\ref{sec:results} the results.
Section~\ref{sec:discussion} discusses limits and Section~\ref{sec:conclusions} concludes.
